\documentclass[12pt]{article}
\usepackage[utf8]{inputenc}
\usepackage[T1]{fontenc}
\usepackage[margin=1in]{geometry}
\usepackage{setspace}
\usepackage{xcolor}
\usepackage[colorlinks=true, urlcolor=blue!60!black, linkcolor=blue!60!black]{hyperref}
\setlength{\parindent}{2em}
\setlength{\parskip}{1.25em}
\renewcommand{\baselinestretch}{1.0}
\begin{document}
\begin{center}
{\LARGE \textit{Nevada Department of Human Resources v. Hibbs}, 538 U.S. 721 (2003)\par}
\vspace{0.8em}
{\large Case Brief --- Four-Page Edition\par}
\end{center}
\vspace{1em}

\section{Memory Jogger}

In a 6--3 decision by Chief Justice Rehnquist, the Supreme Court held that Congress validly abrogated state sovereign immunity under Section 5 of the Fourteenth Amendment in the Family and Medical Leave Act's family-care provision, because sex classifications receive heightened scrutiny and the gender-neutral 12-week leave floor is a congruent and proportional remedy against state sex-role stereotyping (\href{https://supreme.justia.com/cases/federal/us/538/721/}{Justia}).

\section{Detailed Case Facts}

The FMLA of 1993 entitles eligible employees to 12 workweeks of unpaid, job-protected leave a year, including to care for a spouse, child, or parent with a ``serious health condition,'' 29 U.S.C. \S~2612(a)(1)(C). It authorizes damages suits ``against any employer (including a public agency) in any Federal or State court of competent jurisdiction,'' and ``public agency'' includes States. William Hibbs, who worked in the Welfare Division of the Nevada Department of Human Resources, sought leave in 1997 to care for his wife after a serious car accident and neck surgery. The Department approved 12 weeks of intermittent leave. Hibbs stopped working in August; in October the Department told him his leave was exhausted and that he must return by November 12. When he did not, he was fired.

\section{Procedural History}

Hibbs sued the Department and two officers in the District of Nevada for damages and other relief. The district court granted summary judgment for the State on Eleventh Amendment grounds. With the United States intervening, the Ninth Circuit reversed, 273 F.3d 844 (2001), creating a split with the Fifth Circuit's \textit{Kazmier v. Widmann}. The Supreme Court granted certiorari, heard argument January 15, 2003, and affirmed on May 27, 2003.

\section{Judicial Votes}

Affirmed, \textbf{6--3}.
\begin{itemize}
\item \textbf{Majority (5):} Chief Justice Rehnquist, joined by Justices O'Connor, Souter, Ginsburg, and Breyer.
\item \textbf{Concurrences:} Justice Souter, joined by Justices Ginsburg and Breyer; Justice Stevens, in the judgment.
\item \textbf{Dissents:} Justice Scalia; Justice Kennedy, joined by Justices Scalia and Thomas.
\end{itemize}

\section{Holding}

State employees may recover damages in federal court against a nonconsenting State for violating the FMLA's family-care provision. Congress may abrogate immunity if it states its intent unmistakably and acts validly under Section 5; prophylactic legislation must show ``congruence and proportionality between the injury to be prevented or remedied and the means adopted to that end.'' Because sex discrimination triggers heightened scrutiny, it is easier for Congress to show a pattern of unconstitutional state conduct than for age or disability. The holding covers only the family-care provision; \textit{Coleman v. Court of Appeals of Maryland} (2012) later held the self-care provision does not abrogate immunity.

\section{Analysis of Opinions}

\textbf{Majority (Rehnquist).} The FMLA's ``public agency'' language met the clear-statement rule. Since \textit{Seminole Tribe} bars Commerce Clause abrogation, the case turned on Section 5. Rehnquist traced state-sanctioned sex discrimination from \textit{Bradwell} and \textit{Muller} to \textit{Hoyt}, and cited the record: a widening gap between maternity and paternity leave, public-sector policies mirroring private ones, testimony of ``notoriously discriminatory treatment'' of fathers seeking leave, and 15 States offering women extended maternity leave while only four offered men parallel leave. Unlike age and disability in \textit{Kimel} and \textit{Garrett}, sex classifications get heightened scrutiny, so the pattern was easier to show. A mere ban on discriminatory leave would let States offer none. The FMLA's limits---unpaid leave, service requirements, exclusions, notice, and capped damages---made it congruent and proportional.

\textbf{Souter, with Ginsburg and Breyer.} The FMLA was ``undoubtedly valid'' even under the Court's restrictive cases. \textbf{Stevens (judgment).} Doubting the Section 5 rationale, he relied on his view that the Eleventh Amendment's text does not bar a citizen's suit against his own State.

\textbf{Scalia (dissenting).} ``[T]here is no guilt by association'': prophylaxis may extend beyond the violation but not ``beyond the violator,'' so Nevada could demand proof of its own violations.

\textbf{Kennedy, with Scalia and Thomas (dissenting).} The evidence concerned private employers, parenting leave, and federal agencies; 30 States, including Nevada, already had family-care laws. The FMLA is a benefit program, not a remedy.

\section{Examples: Future Applications}

\begin{enumerate}
\item \textbf{Same-Side: Race-Based Disparate Impact.} Congress authorizes damages suits against States for hiring tests that exclude racial minorities without business necessity. Race triggers strict scrutiny, so the evidentiary hurdle is even lower.
\item \textbf{Same-Side: Courthouse Access.} A paraplegic litigant must crawl up courthouse stairs and sues under ADA Title II. \textit{Tennessee v. Lane} (2004) applied \textit{Hibbs}: access to courts is a fundamental right.
\item \textbf{Opposite-Side: FMLA Self-Care Leave.} A state employee fired after seeking leave for his own illness sues for damages. \textit{Coleman} barred the suit; Congress built no record of sex discrimination in sick leave.
\item \textbf{Opposite-Side: Paid Vacation Mandate.} Congress requires four weeks of paid vacation to address wealth inequality. Wealth is not a suspect class, and the benefit exceeds Section 5.
\item \textbf{Fence-Sitter: Pregnancy Accommodation.} A pregnant state worker sues under a reasonable-accommodation statute. Pregnancy lies on the ``faultline between work and family,'' but \textit{Geduldig v. Aiello} treats pregnancy distinctions as not automatically sex-based.
\end{enumerate}

\section{Critique}

\textbf{Originalist and federalism views.} Echoing Kennedy and Scalia, critics argue that \textit{Hibbs} relaxed \textit{Boerne} on thin private-sector and parenting-leave evidence to impose a national labor mandate even on progressive States like Nevada. Some textualists note the irony that the Section 5 gymnastics were needed only because \textit{Hans} and \textit{Seminole Tribe} bar Commerce Clause abrogation.

\textbf{Progressive and doctrinal views.} Robert Post and Reva Siegel read \textit{Hibbs} as vindicating a ``policentric'' model in which Congress enforces evolving equality norms; Samuel Bagenstos praised its ``feminist universalism.'' Erwin Chemerinsky and others object that tying Section 5 power to tiers of scrutiny leaves older and disabled workers unprotected.

\section{Key Quotations}

\begin{quote}
``Stereotypes about women's domestic roles are reinforced by parallel stereotypes presuming a lack of domestic responsibilities for men. Because employers continued to regard the family as the woman's domain, they often denied men similar accommodations or discouraged them from taking leave. These mutually reinforcing stereotypes created a self-fulfilling cycle of discrimination that forced women to continue to assume the role of primary family caregiver, and fostered employers' stereotypical views about women's commitment to work and their value as employees.''
\end{quote}
\noindent --- Chief Justice Rehnquist, Opinion of the Court, 538 U.S. at 736.

\begin{quote}
``By setting a minimum standard of family leave for all eligible employees, irrespective of gender, the FMLA attacks the formerly state-sanctioned stereotype that only women are responsible for family care giving, thereby reducing employers' incentives to engage in discrimination by basing hiring and promotion decisions on stereotypes.''
\end{quote}
\noindent --- Chief Justice Rehnquist, Opinion of the Court, 538 U.S. at 737.

\begin{quote}
``Unlike the statutes at issue in \textit{City of Boerne}, \textit{Kimel}, and \textit{Garrett}, which applied broadly to every aspect of state employers' operations, the FMLA is narrowly targeted at the faultline between work and family---precisely where sex-based overgeneralization has been and remains strongest---and affects only one aspect of the employment relationship.''
\end{quote}
\noindent --- Chief Justice Rehnquist, Opinion of the Court, 538 U.S. at 738.

\end{document}
